% Part: methods
% Chapter: induction
% Section: introduction

\documentclass[../../../include/open-logic-section]{subfiles}

\begin{document}

\olfileid{mth}{ind}{int}

\olsection{प्रस्तावना}

विगमन ही सिद्धतेची महत्त्वाची पद्धती आहे. तर्कशास्त्र, सैद्धान्तिक संगणकशास्त्र आणि गणित यांच्या जवळजवळ सर्व शाखांत ती वेगवेगळ्या रूपांत वापरली जाते. तर्कशास्त्रातील अनेक निष्कर्ष सिद्ध करण्यासाठी ती आवश्यक आहे.

विगमनाचा अनेकदा निगमनाशी विरोध दाखवला जातो आणि विशिष्ट प्रकरणांकडून सामान्य विधानाकडे जाणारे अनुमान, असे त्याचे वर्णन केले जाते. उदाहरणार्थ, अनेक हिरवी पाचूची रत्ने आपण पाहिली आणि पाचू म्हणता येईल असे कोणतेही रत्न हिरवे नसलेले आढळले नाही, तर सर्व पाचू हिरवे असतात असा निष्कर्ष काढू शकतो. हे विगमनाधारित अनुमान आहे: ते अनेक विशिष्ट प्रकरणांकडून---हा पाचू हिरवा आहे, तो पाचू हिरवा आहे, इत्यादी---एका सामान्य विधानाकडे---सर्व पाचू हिरवे असतात---जाते. \emph{गणितीय} विगमनातही सामान्य विधान निष्कर्ष म्हणून मिळते; पण ते निरीक्षणांवरून केलेल्या या “साध्या विगमना”पेक्षा फार वेगळे असते.

स्थूलपणे सांगायचे तर, गणितातील विगमनाधारित सिद्धतेने एखाद्या प्रकारच्या सर्व गणितीय वस्तूंमध्ये विशिष्ट गुणधर्म आहे असा निष्कर्ष मिळतो. सर्वात साध्या प्रकरणात या सिद्धतेतील गणितीय वस्तू नैसर्गिक संख्या असतात. अशा वेळी सर्व नैसर्गिक संख्यांत एखादा गुणधर्म आहे हे सिद्ध करण्यासाठी विगमन वापरतो; त्यात असे दाखवतो:
\begin{enumerate}
\item $0$ मध्ये तो गुणधर्म आहे; आणि
\item एखाद्या~$k$ संख्येत तो गुणधर्म असेल, तर~$k+1$ मध्येही तो असतो.
\end{enumerate}
ज्या गणितीय वस्तूंना संख्या देता येतात त्यांच्याविषयीची सामान्य विधाने सिद्ध करण्यासाठीही नैसर्गिक संख्यांवरील विगमन अनेकदा वापरता येते. उदाहरणार्थ, प्रत्येक सांत संचात सांत संख्येने, समजा~$n$, घटक असतात. विगमन वापरून प्रत्येक~$n$ संख्येसाठी “~$n$ आकारमानाचे सर्व सांत संच \dots{} असतात” हा गुणधर्म आहे असे दाखवता आले, तर सर्व सांत संचांविषयी काहीतरी दाखवले असेल.

\emph{विगमनाधारित व्याख्येने दिलेल्या} गणितीय वस्तूंपर्यंतही विगमनाचा विस्तार करता येतो. उदाहरणार्थ, प्रथम-क्रम तर्कशास्त्रासारख्या औपचारिक भाषेतील अभिव्यक्ती विगमनाधारित व्याख्येने दिलेल्या असतात. \emph{संरचनात्मक विगमन} ही अशा सर्व अभिव्यक्तींविषयी निष्कर्ष सिद्ध करण्याची पद्धती आहे. विशेषतः तर्कशास्त्रात संरचनात्मक विगमन फार उपयुक्त आहे आणि मोठ्या प्रमाणात वापरले जाते.

\end{document}
